\documentclass[11pt]{article}
\usepackage[a4paper,margin=20mm]{geometry}
\usepackage{amsmath,amssymb,xcolor,graphicx}
\usepackage[hypcap=false]{caption}
\usepackage[hidelinks,bookmarksnumbered,bookmarksopen]{hyperref}
\hypersetup{pdftitle={Mechanical properties — Dislocations and strengthening},pdfauthor={Lecturer: Prof. J. Adeyemi},
  pdfcreator={KherveNote}}
\renewcommand\labelitemii{\textbullet}
\renewcommand\labelitemiii{\textbullet}
\renewcommand\labelitemiv{\textbullet}
\renewcommand\labelenumii{\arabic{enumi}.\arabic{enumii}.}
\renewcommand\labelenumiii{\arabic{enumi}.\arabic{enumii}.\arabic{enumiii}.}
\renewcommand\labelenumiv{\arabic{enumi}.\arabic{enumii}.\arabic{enumiii}.\arabic{enumiv}.}
\setlength{\parindent}{0pt}
\setlength{\parskip}{0.6em}
\definecolor{knoteaccent}{HTML}{1A6DD8}
\definecolor{knotemuted}{HTML}{6B6B6B}
\definecolor{knotekey}{HTML}{D96B00}
\newcommand\knotetime[1]{\leavevmode\llap{\color{knotemuted}\scriptsize #1\hspace{1em}}}
\newcommand\knotekey[2][]{\par\noindent#1\colorbox{knotekey!12}{\parbox{\dimexpr\linewidth-2\fboxsep}{%
  \textbf{\color{knotekey}$\star$ Key point.} #2}}\par}
\newcommand\knotequestion[2][]{\par\noindent#1{\color{knoteaccent}\textbf{?}}~\emph{#2}\par}
\newenvironment{knotetranscript}{\par\color{knotemuted}\small}{\par}
\newcommand\knotesummary[1]{\par\noindent\fcolorbox{knoteaccent}{knoteaccent!6}{%
  \parbox{\dimexpr\linewidth-2\fboxsep-2\fboxrule}{\textbf{Summary}\par #1}}\par}
\newenvironment{knotechunk}{}{}
\pagestyle{plain}
\begin{document}
\begin{knotechunk}
{\color{knoteaccent}\rule{\linewidth}{1.5pt}}\par
{\LARGE\bfseries Mechanical properties — Dislocations and strengthening}\par
{\large Lecturer: Prof. J. Adeyemi}\hfill{\color{knotemuted}16 October 2026 \textperiodcentered{} Materials building, LT3}\par
{\color{knoteaccent}\rule{\linewidth}{0.6pt}}\par

\section{Stress and strain}

\knotetime{01:25}Engineering stress $\sigma = F/A_0$, strain $\varepsilon = \Delta L/L_0$; in the elastic range Hooke's law $\sigma = E\varepsilon$.

\begin{itemize}\setlength{\itemsep}{0pt}\setlength{\parskip}{0pt}
  \item yield strength: the 0.2 \% offset stress
  \item ultimate tensile strength (UTS): the maximum of the engineering curve
  \item ductility: elongation at fracture
\end{itemize}

\knotekey[\knotetime{03:05}]{Metals yield at stresses far below the theoretical shear strength (\ensuremath{\approx} G/10) — because of dislocations.}
\end{knotechunk}

\begin{knotechunk}
\section{Dislocations and slip}

\begin{itemize}\setlength{\itemsep}{0pt}\setlength{\parskip}{0pt}
  \item edge dislocation: Burgers vector $\mathbf{b} \perp$ line; screw: $\mathbf{b} \parallel$ line
  \item slip happens on close-packed planes along close-packed directions
  \item FCC: $\{111\}\langle 110\rangle$, 12 slip systems — hence ductile (Cu, Al, Au)
  \item Schmid's law: resolved shear stress $\tau_R = \sigma\cos\phi\cos\lambda$; slip starts when $\tau_R$ reaches the critical value
\end{itemize}
\end{knotechunk}

\begin{knotechunk}
\section{Strengthening mechanisms}

\knotetime{25:25}Every mechanism makes dislocations harder to move:

\begin{enumerate}\setlength{\itemsep}{0pt}\setlength{\parskip}{0pt}
  \item grain refinement — Hall–Petch: $\sigma_y = \sigma_0 + k_y d^{-1/2}$
  \item solid-solution strengthening: solute atoms strain the lattice
  \item precipitation hardening: particles pinned or bypassed by Orowan bowing (Al alloys, Ni superalloys)
  \item work hardening: dislocation density $\rho$ rises, $\tau \approx \alpha G b\sqrt{\rho}$ (Taylor)
\end{enumerate}

\knotekey[\knotetime{27:30}]{Most mechanisms trade strength for ductility; grain refinement improves both, down to the nanocrystalline range.}

\knotequestion[\knotetime{27:55}]{Why does Hall–Petch break down for grain sizes below about 10–20 nm?}
\end{knotechunk}

\begin{knotechunk}
\section{What was said}
\begin{knotetranscript}
\knotetime{10:00:40}Today, why metals are so much weaker than their bonds suggest, and how we make them stronger.

\knotetime{10:07:30}The theoretical shear strength is around G over ten, but real metals yield a thousand times lower.

\knotetime{10:10:50}The answer is dislocations: you move one line of atoms at a time, like a ruck in a carpet.

\knotetime{10:16:20}FCC metals have twelve slip systems, which is why copper and aluminium are so ductile.

\knotetime{10:21:00}Schmid's law: only the shear stress resolved on the slip plane and direction matters.

\knotetime{10:23:50}So every strengthening mechanism is a way of getting in the way of dislocations.

\knotetime{10:28:30}Smaller grains mean more boundaries to pile up against, that's Hall–Petch, one over root d.

\knotetime{10:33:40}Work hardening comes from dislocations tangling with each other, stress goes as root of the density.
\end{knotetranscript}
\end{knotechunk}
\end{document}
